\documentclass[11pt,a4paper]{article}

\usepackage[utf8]{inputenc}
\usepackage[T1]{fontenc}
\usepackage{lmodern}
\usepackage[french]{babel}
\usepackage{amsmath,amssymb,mathtools}
\usepackage{graphicx}
\usepackage{fancyhdr}
\usepackage{array}
\usepackage{enumitem}
\usepackage{xcolor}
\usepackage[hidelinks]{hyperref}

% Marge haute assez grande : les logos font ~2,2 cm
\usepackage[
  a4paper,
  top=3.7cm,
  bottom=2.2cm,
  left=2cm,
  right=2cm,
  headheight=2.6cm,
  headsep=10pt
]{geometry}

\definecolor{mundiblue}{HTML}{0070C0}
\definecolor{mundidark}{HTML}{123B5E}
\definecolor{mundigray}{HTML}{5A5A5A}

\setlength{\parindent}{0pt}
\setlength{\parskip}{6pt}

% En-tête avec deux logos alignés
%---------------------------------------------------
\pagestyle{fancy}
\fancyhf{} % vide l'en-tête et le pied de page

% Les deux logos sont dans la MÊME ligne (tabular*), à hauteur identique.
% \vspace{0pt} aligne le tout en haut de l'en-tête fancyhdr.
\newcommand{\logoentete}[1]{%
  \includegraphics[height=2.2cm,keepaspectratio]{#1}%
}

\fancyhead[C]{%
  \vspace{0pt}%
  \begin{tabular*}{\headwidth}{@{}l@{\extracolsep{\fill}}r@{}}
    \logoentete{honoris.jpg} &
    \logoentete{munidia.png}
  \end{tabular*}%
}

% Numéro de page
\fancyfoot[C]{\thepage}

\renewcommand{\headrulewidth}{0pt} % pas de ligne sous les logos
\setlength{\headheight}{2.6cm}
%---------------------------------------------------

\begin{document}

\begin{center}
{\LARGE\bfseries\color{mundiblue} Université Mundiapolis}\\[0.25cm]
{\large\color{mundigray} Business School}\\[0.15cm]
{\large Licence Management et Gestion des Entreprises}\\[0.45cm]
{\Large\bfseries\color{mundidark} Série n$^{\circ}$ 1}\\[0.2cm]
{\large Chapitre --- Les suites réelles}
\end{center}

\vspace{0.4cm}

\section*{Exercice 1 --- Généralités}

On considère la suite $\bigl(u_n\bigr)$ définie par
\[
u_n = \dfrac{3n+1}{n+2}.
\]

\begin{enumerate}
  \item Calculer $u_0$, $u_1$ et $u_2$.
  \item Calculer $u_{n+1}-u_n$ et en déduire le sens de variation de $\bigl(u_n\bigr)$.
  \item Montrer que, pour tout $n\in\mathbb{N}$, $u_n < 3$. La suite est-elle minorée ?
  \item En déduire que $\bigl(u_n\bigr)$ converge, puis calculer sa limite.
\end{enumerate}

\section*{Exercice 2 --- Calculs de limites}

Déterminer les limites suivantes (justifier, notamment en cas de forme indéterminée) :

\begin{enumerate}
  \item $\displaystyle\lim_{n\to +\infty} \dfrac{4n^3 - n + 2}{2n^3 + 5n - 1}$
  \item $\displaystyle\lim_{n\to +\infty} \bigl(\sqrt{n^2 + n} - n\bigr)$
  \item $\displaystyle\lim_{n\to +\infty} \dfrac{2^n + 3^n}{3^n + 1}$
\end{enumerate}

\section*{Exercice 3 --- Suites arithmétiques et géométriques}

\begin{enumerate}
  \item Une machine-outil est achetée $80\,000$ MAD et amortie linéairement de $8\,000$ MAD par an. On note $V_n$ sa valeur comptable après $n$ années.
  \begin{enumerate}
    \item Justifier que $\bigl(V_n\bigr)$ est une suite arithmétique. Préciser $V_0$ et la raison.
    \item Exprimer $V_n$ en fonction de $n$.
    \item Au bout de combien d'années la machine est-elle totalement amortie ?
  \end{enumerate}
  \item Un capital $C_0 = 15\,000$ MAD est placé à intérêts composés au taux annuel $t=4{,}5\%$. On note $C_n$ le capital acquis après $n$ années.
  \begin{enumerate}
    \item Exprimer $C_n$ en fonction de $n$.
    \item Calculer $C_8$ (valeur exacte, puis arrondie à l'unité).
    \item Au bout de combien d'années le capital aura-t-il doublé ?
  \end{enumerate}
\end{enumerate}

\section*{Exercice 4 --- Suite récurrente}

On considère la suite définie par $u_0 = 1$ et
\[
u_{n+1} = \sqrt{3u_n + 4}.
\]

\begin{enumerate}
  \item Résoudre $\ell = \sqrt{3\ell + 4}$. Quelles sont les valeurs candidates à être limite de $\bigl(u_n\bigr)$ ?
  \item Montrer par récurrence que, pour tout $n\in\mathbb{N}$, $1 \leq u_n \leq 4$.
  \item Montrer que $\bigl(u_n\bigr)$ est croissante.
  \item Conclure.
\end{enumerate}

\end{document}
